\documentclass[10pt,aspectratio=169]{beamer}
\usetheme{Madrid}
\usecolortheme{whale}

\usepackage[utf8]{inputenc}
\usepackage[french]{babel}
\usepackage{graphicx}
\usepackage{booktabs}
\usepackage{listings}
\usepackage{xcolor}
\usepackage{amsmath}

% Configuration de l'affichage du code source
\lstset{
    basicstyle=\tiny\ttfamily,
    stringstyle=\color{blue},
    numbers=none,
    showstringspaces=false,
    breaklines=true,
    frame=single,
    backgroundcolor=\color{gray!5}
}

\title[Présentation du Dataset SETI@home]{Analyse et Exploitation du Dataset Global BOINC / SETI@home (sat09)}
\subtitle{Caractérisation, Standardisation (FTA) et Apport pour la Prédiction en Environnement Contraint}
\author[Mémoire de Master 2]{Présenté par : \textbf{Neussi} \\[0.5em] \small Encadrant : \textbf{Dr. Adamou Hamza}}
\institute[UY1]{Université de Yaoundé I \\ Département d'Informatique \\ Spécialité : Systèmes et Réseaux}
\date[Juin 2026]{Apport et Rôle Scientifique des Traces Globales}

\begin{document}

% Slide de Titre
\begin{frame}
  \titlepage
\end{frame}

% Slide 1: Contexte et Origine des Traces
\begin{frame}{1. Contexte et Origine des Traces (sat09)}
  \begin{block}{Qu'est-ce que SETI@home / BOINC ?}
    \begin{itemize}
      \item \textbf{Calcul Volontaire Pionnier :} Projet de calcul distribué mondial (UC Berkeley) utilisant les ressources de calcul inutilisées (CPU/RAM) des ordinateurs connectés à Internet.
      \item \textbf{Le Failure Trace Archive (FTA) :} Une initiative académique standardisant les traces de disponibilité et de défaillance provenant de divers systèmes distribués mondiaux.
    \end{itemize}
  \end{block}
  
  \begin{columns}[T]
    \begin{column}{0.48\textwidth}
      \begin{alertblock}{La Trace sat09}
        La trace \textbf{sat09} est le jeu de données officiel collecté en instrumentant le client BOINC de SETI@home pour suivre la disponibilité CPU en temps réel.
      \end{alertblock}
    \end{column}
    
    \begin{column}{0.48\textwidth}
      \begin{exampleblock}{Rôle Académique}
        Sert de **baseline de référence mondiale** pour valider les modèles de prédiction de disponibilité dans la littérature scientifique.
      \end{exampleblock}
    \end{column}
  \end{columns}
\end{frame}

% Slide 2: Volumétrie et Métriques
\begin{frame}{2. Volumétrie et Caractéristiques du Dataset}
  Ce dataset se distingue par sa taille massive et sa représentativité globale.
  
  \begin{columns}[T]
    \begin{column}{0.48\textwidth}
      \begin{block}{Chiffres Clés de la Trace}
        \begin{itemize}
          \item \textbf{Nombre de nœuds suivis :} Plus de \textbf{227 210 machines réelles} réparties à travers le monde.
          \item \textbf{Période d'observation :} 1,5 an (de 2007 à 2009).
          \item \textbf{Nature de la mesure :} Disponibilité CPU binaire (1 = CPU disponible pour le calcul, 0 = CPU occupé par l'utilisateur ou machine éteinte).
          \item \textbf{Volume compressé brut :} \textbf{2,18 Go} (BSON/textuel).
        \end{itemize}
      \end{block}
    \end{column}
    
    \begin{column}{0.48\textwidth}
      \begin{exampleblock}{Les Formats Disponibles}
        \begin{itemize}
          \item \textbf{seti09\_raw.tgz (2.18 Go) :} Traces brutes sans aucun filtrage.
          \item \textbf{seti09\_tab.tgz (2.15 Go) :} Format nettoyé en colonnes tabulées (Normalisation standard FTA).
          \item \textbf{seti\_5min\_sql.tgz (2.04 Go) :} Format SQL avec échantillonnage périodique toutes les 5 minutes.
        \end{itemize}
      \end{exampleblock}
    \end{column}
  \end{columns}
\end{frame}

% Slide 3: Pourquoi ce Dataset est qualifié de "Propre" (Clean) ?
\begin{frame}{3. Pourquoi le Dataset est qualifié de ``Propre'' (Clean) ?}
  Le passage de la trace brute (\textit{raw}) à la trace nettoyée (\textit{tabbed}) obéit aux règles strictes du Failure Trace Archive :
  
  \begin{columns}[T]
    \begin{column}{0.48\textwidth}
      \begin{block}{Règles de Filtrage Temporel}
        \begin{itemize}
          \item \textbf{Élimination des micro-intervalles :} Les indisponibilités très courtes ($\le 150$ secondes) survenant après de longues périodes d'activité ($\ge 5$ jours) sont supprimées car jugées non significatives pour l'ordonnancement de tâches.
          \item \textbf{Résolution des conflits temporels :} Correction des anomalies de synchronisation d'horloge entre les machines volontaires.
        \end{itemize}
      \end{block}
    \end{column}
    
    \begin{column}{0.48\textwidth}
      \begin{exampleblock}{Structure Normalisée du Format Tabulaire}
        Chaque ligne du fichier \texttt{seti09\_tab} est formatée selon le standard du FTA :
        \begin{itemize}
          \item \textbf{Event Time :} Date et heure précise du changement d'état.
          \item \textbf{Host ID :} Identifiant unique et anonyme de la machine.
          \item \textbf{State :} Statut de disponibilité (0 ou 1).
        \end{itemize}
      \end{exampleblock}
    \end{column}
  \end{columns}
\end{frame}

% Slide 4: Caractéristiques Comportementales Globales
\begin{frame}{4. Caractéristiques Comportementales Globales}
  Les traces de SETI@home reflètent un contexte d'infrastructure occidentale stable :
  
  \begin{columns}[T]
    \begin{column}{0.48\textwidth}
      \begin{block}{Cycles Circadiens Marqués}
        \begin{itemize}
          \item \textbf{Activité humaine régulière :} Forte corrélation avec l'heure locale (les ordinateurs deviennent indisponibles en journée lorsque les utilisateurs travaillent, et très disponibles la nuit).
          \item \textbf{Saisonnalité hebdomadaire :} Comportement distinct entre les jours de semaine et le week-end.
        \end{itemize}
      \end{block}
    \end{column}
    
    \begin{column}{0.48\textwidth}
      \begin{alertblock}{Stabilité Électrique Absolue}
        \begin{itemize}
          \item \textbf{Aucun délestage :} L'arrêt d'une machine est toujours une décision humaine (extinction volontaire) ou une mise en veille, jamais une coupure de courant accidentelle.
          \item \textbf{Alimentation continue :} Pas de gestion de batterie ou de micro-coupures d'alimentation secteur.
        \end{itemize}
      \end{alertblock}
    \end{column}
  \end{columns}
\end{frame}

% Slide 5: Pourquoi utiliser SETI@home dans notre étude à Yaoundé ?
\begin{frame}{5. Pourquoi utiliser SETI@home dans notre étude à Yaoundé ?}
  Avoir le dataset SETI@home en local est une nécessité méthodologique et scientifique pour notre thèse :
  
  \begin{block}{Les 3 Rôles Clés dans notre Travail}
    \begin{enumerate}
      \item \textbf{Groupe de Contrôle (Baseline comparative) :} Mettre en évidence les spécificités locales de Yaoundé. En comparant les lois de distribution (ex. Weibull), nous prouvons scientifiquement l'effet perturbateur des délestages électriques locaux sur les profils de disponibilité.
      \item \textbf{Pré-entraînement (Transfer Learning) :} Un modèle profond comme le GRU nécessite des millions de lignes pour converger. Nous le pré-entraînons sur SETI@home pour lui apprendre les dynamiques d'activité CPU globales, avant de l'adapter à Yaoundé.
      \item \textbf{Validation de l'Apprentissage Continu (EWC) :} Tester la résilience de nos algorithmes d'adaptation (comme l'Elastic Weight Consolidation) pour s'assurer que le modèle conserve la connaissance globale tout en apprenant les délestages.
    \end{enumerate}
  \end{block}
\end{frame}

% Slide 6: Implémentation du Transfer Learning avec EWC
\begin{frame}{6. Adaptation et Régularisation par EWC}
  Pour éviter que notre GRU compact ne souffre d'**oubli catastrophique** lorsqu'il passe de SETI@home aux traces locales de Yaoundé :
  
  \begin{columns}[T]
    \begin{column}{0.48\textwidth}
      \begin{block}{Le Mécanisme EWC}
        \begin{itemize}
          \item Calcul d'une **matrice d'information de Fisher** ($F$) sur le dataset SETI@home pour identifier les poids critiques du modèle.
          \item Ajout d'une pénalité quadratique lors de l'entraînement sur les données locales de Yaoundé.
        \end{itemize}
      \end{block}
    \end{column}
    
    \begin{column}{0.48\textwidth}
      \begin{exampleblock}{Formulation Mathématique}
        $$\mathcal{L}(\theta) = \mathcal{L}_{\text{Yaoundé}}(\theta) + \sum_{i} \frac{\lambda_{\text{EWC}}}{2} F_i (\theta_i - \theta_{\text{SETI}, i})^2$$
        \begin{itemize}
          \item $\theta_i$ : Paramètres actuels du modèle.
          \item $\theta_{\text{SETI}, i}$ : Poids optimaux appris sur SETI@home.
          \item $\lambda_{\text{EWC}}$ : Force de rétention de la mémoire globale.
        \end{itemize}
      \end{exampleblock}
    \end{column}
  \end{columns}
\end{frame}

% Slide 7: RLS-ARX : Complémentarité Frugale (TinyML)
\begin{frame}{7. RLS-ARX : Complémentarité Frugale (TinyML)}
  En parallèle du GRU profond, le modèle linéaire **RLS-ARX** gère l'adaptation continue locale :
  
  \begin{columns}[T]
    \begin{column}{0.48\textwidth}
      \begin{block}{Frugalité Computationnelle}
        \begin{itemize}
          \item \textbf{Zéro stockage d'historique :} Les poids et la covariance inverse ($P_t$) sont mis à jour récursivement sans conserver de données en RAM ($< 1$ Ko).
          \item \textbf{Mémoire temporelle effective :} Configurée via le facteur d'oubli $\lambda = 0.99977$ pour équivaloir à une fenêtre glissante de 72 heures.
        \end{itemize}
      \end{block}
    \end{column}
    
    \begin{column}{0.48\textwidth}
      \begin{exampleblock}{Immunité contre l'Oubli}
        \begin{itemize}
          \item \textbf{Convexité stricte :} Résolution d'un problème quadratique convexe garantissant un optimum global unique.
          \item \textbf{Covariance comme régularisateur :} La matrice de covariance inverse $P_t$ protège les directions caractéristiques stables contre le bruit transitoire.
        \end{itemize}
      \end{exampleblock}
    \end{column}
  \end{columns}
\end{frame}

% Slide 8: Méthodologie d'Exploitation Frugale en Python
\begin{frame}[fragile]{8. Méthodologie d'Exploitation Frugale (Python)}
  Pour manipuler les 2,15 Go de traces sans saturer la mémoire vive lors des phases de pré-entraînement :
  
  \begin{lstlisting}[language=python, basicstyle=\tiny\ttfamily, stringstyle=\color{green!50!black}, keywordstyle=\color{blue}]
# Lecture iterative par blocs (streaming) des traces tabulaires
def stream_seti_traces(file_path):
    with open(file_path, "r") as f:
        for line in f:
            if line.startswith("#") or not line.strip():
                continue
            # Parsing d'une ligne normalisee FTA
            parts = line.strip().split("\t")
            event_time = float(parts[0])
            host_id = int(parts[1])
            state = int(parts[2])
            
            yield event_time, host_id, state

# Exemple d'usage pour calculer le taux de disponibilite global
total_events = 0
active_events = 0
for t, host, state in stream_seti_traces("research_models/data_boinc/raw/seti09_tab.txt"):
    total_events += 1
    if state == 1:
        active_events += 1
print(f"Disponibilite globale : {active_events / total_events * 100:.2f}%")
  \end{lstlisting}
\end{frame}

% Slide 9: Conclusion et Perspectives
\begin{frame}{9. Conclusion et Perspectives}
  \begin{block}{Un Couplage de Données Robuste et Rigoureux}
    L'intégration locale du dataset SETI@home aux côtés de nos données réelles de Yaoundé (VC-UY1) apporte à la recherche :
    \begin{itemize}
      \item \textbf{Solidité comparative :} Une preuve claire et formelle de l'impact de l'instabilité électrique locale sur la disponibilité système.
      \item \textbf{Rigueur méthodologique :} Une validation du transfert de connaissances (Transfer Learning) d'une échelle mondiale à une échelle locale contrainte.
      \item \textbf{Pertinence TinyML :} Une architecture hybride (RLS-ARX frugal et GRU-EWC quantifié) optimisant la précision de prédiction sous contrainte stricte de RAM ($< 10$ Mo).
    \end{itemize}
  \end{block}
  
  \centering
  \vspace{1em}
  \Large\textbf{Merci pour votre attention. Place aux questions.}
\end{frame}

\end{document}
